\scanpage{10}\mspage{9}
\makeatletter\tagsleft@true\makeatother
de \uline{compatibilité} dans $\unsure{\mathcal{M}|M}$ est une \deleted{notion} relation importante, qui se \unsure{surajoute|surrajoute} aux relations $\underset{\mathcal{M}}{\leq}$, $\underset{\mathcal{M}}{\ll}$. Moralement, pour récupérer l'atelier (à iso unique près) il suffit de connaitre sur $\mathcal{M}$ ces deux relations d'ordre, et la relation de compatibilité.
% slanted
\margin{au moins dans le cas modéré}

Les relations \deleted{\unsure{suivantes}} \inserted{(compatibilité, disjonction)} sont si importantes, qu'elles méritent chacune une notation lapidaire. Je propose
\[
\deleted{F\parallel G}\qquad F\parallel G \tag*{(1.14)}
\]
pour la disjonction, et
\[
\deleted{F\unsure{\mathrel{\tikz[baseline=0pt,x=1em,y=1em]{\draw (-.38,.62)--(.38,.25)--(-.43,-.10); \draw (.34,.62)--(-.32,.25)--(.38,-.10);}}|\mathrel{\tikz[baseline=0pt,x=1em,y=1em]{\draw (-.44,.26)--(0,.58)--(.44,.26)--(0,-.06)--cycle; \draw (0,.58)--(0,-.06);}}}G}\qquad F\mathrel{\tikz[baseline=0pt,x=1em,y=1em]{\draw (-.44,.26)--(0,.58)--(.44,.26)--(0,-.06)--cycle; \draw (0,.58)--(0,-.06);}} G\quad\text{ou}\quad
\tikz[baseline=(fg.base)]{\node[inner sep=1pt] (fg) {$F\mathrel{\tikz[baseline=0pt,x=1em,y=1em]{\draw (-.38,.62)--(.38,.25)--(-.43,-.10); \draw (.34,.62)--(-.32,.25)--(.38,-.10);}} G$}; \draw ([shift={(70:2.12em and 1.40em)}]fg.center) arc[start angle=70,end angle=410,x radius=2.12em,y radius=1.40em];}
\tag*{(1.15)}
\]
pour la compatibilité (NB $\mathrel{\tikz[baseline=0pt,x=1em,y=1em]{\draw (-.38,.62)--(.38,.25)--(-.43,-.10); \draw (.34,.62)--(-.32,.25)--(.38,-.10);}}$ est plus rapide \unsure{:|à} écrire que $\mathrel{\tikz[baseline=0pt,x=1em,y=1em]{\draw (-.44,.26)--(0,.58)--(.44,.26)--(0,-.06)--cycle; \draw (0,.58)--(0,-.06);}}$, $\unsure{\mathrel{\tikz[baseline=0pt,x=1em,y=1em]{\draw (-.12,.56)--(-.48,.25)--(-.12,-.02); \draw (.10,.56)--(.49,.25)--(.13,-.03); \draw (-.32,.51)--(-.59,.25)--(-.28,.18)--(.20,.30)--(.49,.25); \draw (-.28,.18)--(.02,.45); \draw (.26,.55)--(.50,.25);}}|\deleted{\mathrel{\tikz[baseline=0pt,x=1em,y=1em]{\draw (-.12,.56)--(-.48,.25)--(-.12,-.02); \draw (.10,.56)--(.49,.25)--(.13,-.03); \draw (-.32,.51)--(-.59,.25)--(-.28,.18)--(.20,.30)--(.49,.25); \draw (-.28,.18)--(.02,.45); \draw (.26,.55)--(.50,.25);}}}}$ mais est moins suggestif de l'idée combinatoire exprimée.

On n'a rien dit pour l'instant au sujet de la relation $\deleted{\ll}\preccurlyeq$ de raffinement sur $\unsure{\mathcal{F}|\mathfrak{F}}$. Elle est formellement de nature plus délicate, et est la notion la plus spécifique dans l'``Algèbre des figures''. La signification géométrique s'éclaire par les \unsure{exemples|exemple}, et les axiomes que \unsure{ceux-|ceux}%
